\chapter{The ErisML language and its intermediate representation}\label{ch:erisml}

A scene in this monograph is one file. It names the people in a situation, the promises that bind
them, the norms that govern an agent, the events that can happen and the regions of moral space the
situation can occupy. This chapter describes that file and the typed record the compiler loads it
into. It covers the two textual forms that carry the name ErisML, the loader, the schema of the
intermediate representation, norms and the condition tokens that say when a norm is in force, the
declarations a scene keeps in its \code{extra} block, and the typed graph whose hash enters the
audit chain. The running example is Margaret's home, the scene that drives the robot of
Part~III (Chapter~\ref{ch:home}). Chapter~\ref{ch:runtime} describes how the scene runtime executes
what this chapter declares.
% src: erisml-compiler src/erisml_compiler/ingestion/structured_loader.py:1-5
% src: gtc-prototype twin/scene/margaret_home.erisml:1-25

\section{Two textual forms carry the name ErisML}\label{sec:erisml-two-forms}

The first form is a grammar in the erisml-lib repository. It is written for the Lark parser and is
headed ErisML Grammar v2.0 with support for \deme{} 2.0. A model has one environment block, one or
more agent blocks, and optional norms, ethics and dynamics blocks. The environment declares object
types and typed state variables. An agent declares capabilities, beliefs, intents and constraints.
The norms block holds prohibition, obligation and sanction rules. Each rule is an expression over
names, numbers and strings, built with comparison and Boolean operators.
% src: erisml-lib src/erisml/language/grammar.lark:1-6 (header, model rule)
% src: erisml-lib src/erisml/language/grammar.lark:11-43 (environment, agent, norms)
% src: erisml-lib src/erisml/language/grammar.lark:93-105 (expressions)

\begin{lstlisting}[caption={The norms block of the erisml-lib grammar}, label={lst:lark-norms}]
norms_block: "norms" CNAME "{" norm_rule* "}"
norm_rule: prohibition_rule
         | obligation_rule
         | sanction_rule

prohibition_rule: "prohibition:" expr ";"
obligation_rule: "obligation:" expr ";"
sanction_rule: "sanction:" expr ";"
\end{lstlisting}
% src: erisml-lib src/erisml/language/grammar.lark:36-43

The ethics block declares weights on six MoralVector dimensions, the configuration of ethics-module
tiers 0 to 4, module declarations with a tier, a weight, a veto flag and tags, and veto rules that
fire when an expression holds. The dynamics block declares joint-action costs and reward terms.
% src: erisml-lib src/erisml/language/grammar.lark:48-88

The parser builds an LALR table from this grammar and transforms the tree into a typed abstract
syntax tree. Its final step collects the environment, the agents and the norms, and it rejects a
model without an environment. It does not attach an ethics declaration to the tree, although the
tree's schema has a field for one.
% src: erisml-lib src/erisml/language/parser.py:43 (Lark(..., parser="lalr"))
% src: erisml-lib src/erisml/language/parser.py:171-184 (model collects env, agents, norms only)
% src: erisml-lib src/erisml/language/ast.py:113-132 (EthicsDecl, ModelAST.ethics)

The second form is the one the twin uses. The compiler calls it ErisML source and ties it to
requirement FR-19 of its design specification. It is YAML whose top-level keys are the field names
of the compiler's intermediate representation. The compiler's code generator writes this form by
dumping the representation as YAML. Its own docstring calls the format minimal and YAML-like and
says that a later phase will replace it with the formal grammar.
% src: erisml-compiler src/erisml_compiler/ingestion/structured_loader.py:1-5
% src: erisml-compiler src/erisml_compiler/erisml_backend/codegen.py:1-6,15-38

The two forms are not yet joined. No module under the compiler's source tree imports the erisml-lib
language package or its grammar. In the rest of this monograph ErisML means the second form unless
a sentence says otherwise.
% src: grep for "erisml.language|parse_erisml|grammar.lark" over erisml-compiler src/ returned nothing (2026-10-03)

\begin{figure}[t]
\centering
% Chapter 4. From a scene file to a validated record, a typed graph and its hash.
% src: erisml-compiler src/erisml_compiler/ingestion/structured_loader.py:45-76
% src: erisml-compiler src/erisml_compiler/ir/schemas.py:394-458
% src: erisml-compiler docs/architecture.md:106-127, src/erisml_compiler/ir/graph/canonical.py:48-82
\begin{tikzpicture}[house, node distance=4mm]
\def\pw{25mm}
\def\ph{47mm}
% five numbered panels
\foreach \i/\c/\x in {1/Purple/0, 2/Blue/33, 3/Green/66, 4/Orange/99, 5/Teal/132} {
  \node[panel=\c, minimum width=\pw+3mm, minimum height=\ph, anchor=north] (p\i) at (\x mm, 0) {};
  \node[step=\c] (s\i) at ([yshift=-5.5mm]p\i.north) {\i};
}
\node[stepname=Purple, below=1.2mm of s1, text width=\pw, align=center] (n1) {Scene source};
\node[desc, below=0.3mm of n1, text width=\pw] (d1) {ErisML source,\\YAML};
\node[detail, below=1.5mm of d1, text width=\pw-2mm] {\texttt{margaret\_home}\\\texttt{.erisml}\\IR field names\\as top-level keys\\an \texttt{extra} block};

\node[stepname=Blue, below=1.2mm of s2, text width=\pw, align=center] (n2) {Loader};
\node[desc, below=0.3mm of n2, text width=\pw] (d2) {\texttt{load\_structured}\\\texttt{\_input}};
\node[detail, below=1.5mm of d2, text width=\pw-2mm] {.erisml .yaml .yml\\read as YAML\\SHA-256 of the text\\each section\\validated};

\node[stepname=Green, below=1.2mm of s3, text width=\pw, align=center] (n3) {CompilerIR};
\node[desc, below=0.3mm of n3, text width=\pw] (d3) {schema\\\texttt{ir\_v0.3}};
\node[detail, below=1.5mm of d3, text width=\pw-2mm] {stakeholders\\events\\commitments\\norms\\\texttt{extra}};

\node[stepname=Orange, below=1.2mm of s4, text width=\pw, align=center] (n4) {MoralGraph};
\node[desc, below=0.3mm of n4, text width=\pw] (d4) {Pass 7.5\\\texttt{graph\_from\_flat}};
\node[detail, below=1.5mm of d4, text width=\pw-2mm] {6 node kinds\\11 edge kinds\\payloads carry\\the records\\verbatim};

\node[stepname=Teal, below=1.2mm of s5, text width=\pw, align=center] (n5) {Graph hash};
\node[desc, below=0.3mm of n5, text width=\pw] (d5) {canonical JSON\\then SHA-256};
\node[detail, below=1.5mm of d5, text width=\pw-2mm] {nodes by id\\edges by source,\\target, kind,\\payload\\\texttt{graph\_hash}};

\foreach \a/\b in {1/2, 2/3, 3/4, 4/5} { \draw[flow] (p\a.east) -- (p\b.west); }
% the runtime branch
\node[panel=Grey, below=7mm of p3, text width=52mm] (rt) {\textcolor{hsGrey}{\bfseries Scene runtime}\\[1pt]{\scriptsize\color{hsTextMid} reads norms and \texttt{extra}, Chapter 7}};
\draw[flow] (p3.south) -- node[flowlabel, right]{the same record} (rt.north);
\draw[flow] (p4.north) .. controls +(0,9mm) and +(0,9mm) .. node[flowlabel, above]{back to flat lists, \texttt{flat\_from\_graph}} (p3.north);
\end{tikzpicture}

\caption{The path from a scene file to a validated record, a typed graph and its hash. The scene
runtime reads the validated record directly and needs none of the later passes.}
\label{fig:erisml-source-to-hash}
\end{figure}

\section{The loader turns a file into a validated record}\label{sec:erisml-loader}

The function \code{load\_\allowbreak{}structured\_\allowbreak{}input} reads a scene from disk
(Figure~\ref{fig:erisml-source-to-hash}). A file whose name ends in \code{.erisml}, \code{.yaml} or
\code{.yml} is parsed as YAML with a safe loader. Any other file is parsed as JSON. The loader
computes the SHA-256 digest of the raw text and stores it in the document record, together with the
path and the time of loading. The document's identifier and title come from a \code{document}
block when the file has one, and from the file name otherwise. The original Tier-1 input format
named its parties \code{agents}, and the loader still accepts that key in place of
\code{stakeholders}.
% src: erisml-compiler src/erisml_compiler/ingestion/structured_loader.py:45-66

Each section is then built into its Pydantic model, so a field of the wrong type or a value outside
a declared set stops the load. The loader fills the document, the stakeholders, the relations, the
events, the commitments, the norms, the ethical facts and the \code{extra} block. The compiler's
pipeline fills the remaining fields in later passes, which Chapters~5 and~6 describe. The scene
runtime needs only what the loader fills. The twin builds its runtime from the loader's output with
no pipeline pass in between.
% src: erisml-compiler src/erisml_compiler/ingestion/structured_loader.py:20-22 (partially populated IR), 67-76
% src: gtc-prototype twin/brain.py:191-194, 312-320 (SceneRuntime(load_structured_input(...)))

The document's \code{raw\_text} matters beyond bookkeeping. It is the scene's description in plain
words, and both models of the scene agent receive it with every call
(Section~\ref{sec:runtime-agent}). For Margaret's home it says that she lives alone with her dog,
that she has refused recording and sharing of her data, that her bedroom is private, and that visits
are arranged through the monitoring centre.
% src: erisml-compiler src/erisml_compiler/runtime/agent.py:156-161, 245-248 ("scene": raw_text)
% src: gtc-prototype twin/scene/margaret_home.erisml:33-40

\section{The intermediate representation is a typed record}\label{sec:erisml-ir}

The intermediate representation is the Pydantic model \code{CompilerIR}. Its schema version is
\code{erisml\_\allowbreak{}compiler\_\allowbreak{}ir\_\allowbreak{}v0.3}. Assignments to its fields are validated as well as its
construction. Figure~\ref{fig:erisml-ir-anatomy} groups its fields by who fills them.
% src: erisml-compiler src/erisml_compiler/ir/schemas.py:394-458
% src: erisml-compiler src/erisml_compiler/__init__.py:12 (__schema_version__)

\begin{figure}[t]
\centering
% Chapter 4. The fields of CompilerIR, grouped by who fills them, and the extra keys the runtime reads.
% src: erisml-compiler src/erisml_compiler/ir/schemas.py:394-458
% src: erisml-compiler src/erisml_compiler/ingestion/structured_loader.py:54-76
% src: erisml-compiler src/erisml_compiler/runtime/scene.py:11-25,109-129; runtime/agent.py:5-21,105-113
% src: erisml-compiler src/erisml_compiler/process/model.py:139-251
\begin{tikzpicture}[house]
\node[panel=Green, text width=44mm, minimum height=44mm, anchor=north] (load) at (0,0) {};
\node[step=Green] (s1) at ([yshift=-5.5mm]load.north) {1};
\node[stepname=Green, below=1mm of s1] (n1) {Filled by the loader};
\node[desc, below=0.3mm of n1, text width=42mm] (d1) {what a scene file declares};
\node[detail, below=1.5mm of d1, text width=40mm] {\texttt{document}\\\texttt{stakeholders}\quad\texttt{relations}\\\texttt{events}\quad\texttt{commitments}\\\texttt{norms}\quad\texttt{ethical\_facts}\\\texttt{extra}};

\node[panel=Grey, text width=44mm, minimum height=44mm, anchor=north] (pipe) at (52mm,0) {};
\node[step=Grey] (s2) at ([yshift=-5.5mm]pipe.north) {2};
\node[stepname=Grey, below=1mm of s2] (n2) {Filled by the pipeline};
\node[desc, below=0.3mm of n2, text width=42mm] (d2) {Chapters 5 and 6};
\node[detail, below=1.5mm of d2, text width=40mm] {\texttt{segments}\quad\texttt{conflicts}\\\texttt{moral\_vectors}\quad\texttt{timeline}\\\texttt{moral\_tensor\_v3}\\\texttt{em\_outputs}\quad\texttt{deme\_verdict}\\\texttt{graph}\quad\texttt{projections}\\\texttt{decision\_proof}\quad\texttt{audit}};

\node[panel=Amber, text width=44mm, minimum height=44mm, anchor=north] (extra) at (104mm,0) {};
\node[step=Amber] (s3) at ([yshift=-5.5mm]extra.north) {3};
\node[stepname=Amber, below=1mm of s3] (n3) {Keys of \texttt{extra}};
\node[desc, below=0.3mm of n3, text width=42mm] (d3) {read by the runtime, the agent\\and the process loader};
\node[detail, below=1.5mm of d3, text width=40mm] {\texttt{agent}\quad\texttt{role}\\\texttt{capabilities}\quad\texttt{conditions}\\\texttt{oversight}\quad\texttt{strata}\\\texttt{event\_types}\quad\texttt{actors}\\\texttt{default\_action}\\\texttt{processes}};

\draw[flow] (load.east) -- (pipe.west);
\draw[flow] (load.north) .. controls +(0,11mm) and +(0,11mm) .. node[flowlabel, above]{the \texttt{extra} block of panel 1, opened up} (extra.north);
\end{tikzpicture}

\caption{The fields of the intermediate representation grouped by who fills them. Everything the
runtime reads sits in panel 1, and most of it inside \code{extra}.}
\label{fig:erisml-ir-anatomy}
\end{figure}

\paragraph{Stakeholders and relations.}
A stakeholder has an identifier, a label, a type from eight values and a list of roles from thirteen.
The roles include agent, patient, beneficiary, authority and protector. Three further fields are
free strings by convention, namely agency, vulnerability and consent status. A relation is a typed
link from one stakeholder to another. Margaret's home declares four stakeholders and two relations.
The robot cares for Margaret, and the monitoring centre oversees the robot.
Listing~\ref{lst:scene-margaret} shows Margaret's record.
% src: erisml-compiler src/erisml_compiler/ir/schemas.py:83-123
% src: gtc-prototype twin/scene/margaret_home.erisml:41-80

\begin{lstlisting}[caption={Margaret as a stakeholder (\code{margaret\_home.erisml}, lines 42 to 48)}, label={lst:scene-margaret}]
- id: margaret
  label: Margaret, who lives alone
  type: individual
  roles: [patient, beneficiary]
  agency: full
  vulnerability: high
  consent_status: not_obtained
\end{lstlisting}
% src: gtc-prototype twin/scene/margaret_home.erisml:42-48

Two of these fields create runtime state. A stakeholder whose consent status is set and is not
\code{n/a} receives a consent machine, and a stakeholder in the authority or protector role receives
a legitimacy machine (Section~\ref{sec:runtime-machines}).
% src: erisml-compiler src/erisml_compiler/runtime/scene.py:119-128

\paragraph{Events.}
An event has an identifier, a time index that may not be negative, a type, and optional actor,
target, content and conditions. The type is a plain string in the schema. The vocabulary of types
belongs to the scene and is declared in \code{extra} (Section~\ref{sec:erisml-event-types}). A scene
file may list standing events. Margaret's home lists four, at time indices 0 to 3. Margaret refuses
consent to recording, data sharing and entry to her bedroom. The centre assigns the robot its role.
The centre delegates elevated action to the robot only on a governor ruling of elevate. Margaret
starts reading.
% src: erisml-compiler src/erisml_compiler/ir/schemas.py:126-134
% src: gtc-prototype twin/scene/margaret_home.erisml:110-134

\paragraph{Commitments.}
A commitment has a type from seven (vow, promise, contract, oath, covenant, role duty, policy), a
holder, a beneficiary, a content string, a voluntariness, a status from seven and a legitimacy from
six. Its \code{defeasibility\_\allowbreak{}conditions} name conditions under which it may be overridden. The robot
holds three commitments to Margaret. The privacy promise is the one with a defeasibility condition
(Listing~\ref{lst:scene-privacy}).
% src: erisml-compiler src/erisml_compiler/ir/schemas.py:142-163
% src: gtc-prototype twin/scene/margaret_home.erisml:81-109

\begin{lstlisting}[caption={The privacy promise (lines 91 to 100)}, label={lst:scene-privacy}]
- id: privacy_promise
  type: promise
  holder: robot
  beneficiary: margaret
  content: do not record Margaret, share her data, or enter her bedroom without her permission
  voluntariness: voluntary
  created_at_event: e0
  status: active
  legitimacy: fully_legitimate
  defeasibility_conditions: [corroborated_life_threatening_emergency]
\end{lstlisting}
% src: gtc-prototype twin/scene/margaret_home.erisml:91-100

\paragraph{Ethical facts and the remaining fields.}
An ethical fact has a kind from fourteen, subjects, a description, a severity and a confidence.
Margaret's home declares three, about her vulnerability, her refusal of consent, and the conflict
between duty of care and privacy in an emergency. The fields that hold moral vectors, tensors,
module outputs, verdicts, projections, the decision proof and the audit record are filled by the
pipeline.
% src: erisml-compiler src/erisml_compiler/ir/schemas.py:197-222, 405-458
% src: gtc-prototype twin/scene/margaret_home.erisml:135-150

\section{A norm binds an actor to an action under conditions}\label{sec:erisml-norms}

Besides an identifier and source spans, a norm record has eight fields. Its \code{modality} is one of prohibition, obligation, permission and
exception. Its \code{actor} names who it binds, its \code{action} names what, and its optional
\code{target} names toward whom. Its \code{priority\_tier} is a non-negative integer with default 1.
A lower tier is more urgent, and tier 0 is the most urgent. Its \code{defeasible} flag defaults to
true. Its \code{source} is a free string. Its \code{conditions} list the condition tokens that must
all hold for the norm to be in force, and an empty list means always in force.
% src: erisml-compiler src/erisml_compiler/ir/schemas.py:166-189
% src: erisml-compiler src/erisml_compiler/runtime/scene.py:276 (lower tier is higher priority)
% src: erisml-compiler src/erisml_compiler/runtime/agent.py:44-45 (tier 0 the most urgent)

The source string carries one convention the runtime reads. A defeasible norm whose source is
\code{commitment:}\textit{id} is lifted while that commitment's machine is in any state other than
active. Margaret's three privacy prohibitions (record, share data, enter the bedroom) have the source
\code{commitment:privacy\_\allowbreak{}promise}, so they are lifted when the promise is overridden.
% src: erisml-compiler src/erisml_compiler/runtime/scene.py:14-15, 241-249
% src: gtc-prototype twin/scene/margaret_home.erisml:152-176

Listing~\ref{lst:scene-ems} shows two norms on the same action. The prohibition is non-defeasible,
rests on the delegation event \code{e2}, and stands unless one of two named conditions holds. The
obligation rests on the duty of care and is in force exactly when the first condition holds. Both
are at tier 0.

\begin{lstlisting}[caption={Two norms on calling emergency services (lines 187 to 208)}, label={lst:scene-ems}]
# the robot calls emergency services only on a governor ruling: elevate (a corroborated
# emergency) or authorize_ems (the centre unreachable and one attested physical witness). No words,
# screen, message or unattested device can send emergency services to her door.
- id: n6p
  modality: prohibition
  actor: robot
  action: call_emergency_services
  target: emergency_services
  priority_tier: 0
  defeasible: false
  source: delegation_e2
  conditions: ["not:cond:corroborated_life_threatening_emergency", "not:cond:ems_authorized"]
# ---- a corroborated emergency: emergency services at once
- id: n6
  modality: obligation
  actor: robot
  action: call_emergency_services
  target: emergency_services
  priority_tier: 0
  defeasible: false
  source: commitment:duty_of_care
  conditions: ["cond:corroborated_life_threatening_emergency"]
\end{lstlisting}
% src: gtc-prototype twin/scene/margaret_home.erisml:187-208

A norm's condition list is a conjunction. A scene expresses a disjunction by writing several norms
with the same modality and action. Six obligations in Margaret's home oblige a check-in, each on a
different event (inactivity, an unconfirmed alarm, a fall, inflicted harm, medical distress, a danger
in the home). The action is obliged when any of them is in force.
% src: gtc-prototype twin/scene/margaret_home.erisml:229-283 (n4, n4b, n4c, n4d, n4e, n4f)
% src: erisml-compiler src/erisml_compiler/runtime/scene.py:283-287 (any obligation in force)

One action name is reserved. A norm whose action is \code{elevated\_\allowbreak{}action} applies to every
capability the scene marks as elevated. Margaret's norm \code{n3} prohibits all seven elevated
capabilities at tier 0 outside a corroborated emergency.
% src: erisml-compiler src/erisml_compiler/runtime/scene.py:18-19, 67, 253-260
% src: gtc-prototype twin/scene/margaret_home.erisml:177-186, 1152-1177

Figure~\ref{fig:erisml-census} counts what the scene declares. Of its 65 norms, 45 are obligations,
19 are prohibitions and one is a permission. The permission allows chores at tier 3. The scene
declares no exception.
% src: gtc-prototype twin/scene/margaret_home.erisml:151-770 (counted with PyYAML, 2026-10-03)
% src: gtc-prototype twin/scene/margaret_home.erisml:763-770 (the one permission, chores, tier 3)

\begin{figure}[t]
\centering
% Chapter 4. What Margaret's scene declares, counted from the file (twin/scene/margaret_home.erisml).
% Counts taken 2026-10-03 by loading the file with PyYAML and counting list entries:
% norms 65 (45 obligation, 19 prohibition, 1 permission); events 4; stakeholders 4; commitments 3;
% event_types 48 (24 with source: system); capabilities 25 (7 elevated, 4 governed);
% strata 7; extra.conditions 12; processes 1 (escalation_ladder, 23 nodes, 23 flows).
\begin{tikzpicture}[house]
\def\tw{33mm}
\foreach \c/\x/\y/\num/\what/\more in {
  Purple/0/0/4/stakeholders/{Margaret, the robot,\\the centre,\\emergency services},
  Purple/40/0/3/commitments/{duty of care,\\privacy promise,\\household chores},
  Blue/80/0/65/norms/{45 obligations,\\19 prohibitions,\\1 permission},
  Blue/120/0/4/standing events/{consent refused,\\role assigned,\\authority delegated,\\an activity started},
  Teal/0/-33/48/event types/{24 of them from\\the system},
  Green/40/-33/25/capabilities/{7 elevated,\\4 governed},
  Teal/80/-33/7/stratifications/{12 named\\conditions},
  Orange/120/-33/1/process/{the escalation ladder,\\23 nodes, 23 flows}} {
  \node[panel=\c, text width=\tw, minimum height=30mm, anchor=north] (p) at (\x mm, \y mm) {};
  \node[font=\sffamily\bfseries\Large, text=hs\c, below=1.5mm of p.north] (n) {\num};
  \node[stepname=\c, below=0mm of n] (w) {\what};
  \node[detail, below=1mm of w, text width=\tw-2mm] {\more};
}
\end{tikzpicture}

\caption{What the scene of Margaret's home declares, counted from the file. Half of the 48 event
types come only from the system, never from a model's reading of the world.}
\label{fig:erisml-census}
\end{figure}

The deontic modalities themselves are standard in normative multi-agent systems
\cite{boella2006normative}. The SLEEC line of work defines a language for social, legal, ethical,
empathetic and cultural rules for autonomous agents, with defeaters, conflict and redundancy checks,
and verification against robot models, and it uses assistive and care robots as running examples
\cite{getiryaman2023sleec,getiryaman2024sleectoolkit,getiryaman2025sleecjss,ribeiro2026sleecframework}.
Feng et al.\ use language models to help operationalize such rules \cite{feng2024normative}, and
De~Sanctis et al.\ enforce them at runtime on an assistive robot \cite{desanctis2026sleecruntime}.
Joshi et al.\ apply deontic policies to the runtime governance of agentic systems
\cite{joshi2026deontic}. ErisML's norm record does not claim to improve on these languages.
Chapter~\ref{ch:runtime} and Part~III describe what the scene adds around it, which is the control
of who may assert the events that put a norm in force.
% src: gtc-prototype paper/aies2027/prior_art.md:67-83 (credit SLEEC prominently, no novelty for the norm language)

\section{Condition tokens form a small conjunctive language}\label{sec:erisml-tokens}

A condition token is a short string with a kind prefix. Six kinds exist
(Figure~\ref{fig:erisml-tokens}). The runtime evaluates them in the method \code{holds}, and the
process loader classifies them in \code{TokenChecker}.
% src: erisml-compiler src/erisml_compiler/runtime/scene.py:151-189
% src: erisml-compiler src/erisml_compiler/process/model.py:94-136

\begin{figure}[t]
\centering
% Chapter 4. The six condition token forms, what each reads, and when the process loader calls it structural.
% src: erisml-compiler src/erisml_compiler/runtime/scene.py:151-189 (holds)
% src: erisml-compiler src/erisml_compiler/process/model.py:94-136 (TokenChecker.check)
\begin{tikzpicture}[house]
\def\tw{43mm}
\foreach \i/\c/\x/\y/\tok/\reads/\st in {
  1/Blue/0/0/{event:\textit{type}[=\textit{content}]}/{an event of that type, and content,\\since the discharge point}/{structural when the type\\is a system event},
  2/Blue/52/0/{latest:\textit{type}=\textit{content}}/{the newest event of that type\\since then has that content}/{structural when the type\\is a system event},
  3/Purple/104/0/{cond:\textit{name}}/{every token of a condition\\named in \texttt{extra.conditions}}/{structural when all\\of its tokens are},
  4/Green/0/-35/{state:\textit{machine}=\textit{state}}/{a commitment, consent, legitimacy\\or stratum machine is in that state}/{always structural},
  5/Teal/52/-35/{stratum:\textit{name}=\textit{state}}/{the stratification is in that state,\\entered since the discharge point}/{structural when every gate into\\or out of the state fires\\on a system event},
  6/Grey/104/-35/{not:\textit{token}}/{the negation, with a stratum\\read as it stands now}/{structural when the token\\it negates is}} {
  \node[panel=\c, text width=\tw, minimum height=31mm, anchor=north] (p\i) at (\x mm, \y mm) {};
  \node[step=\c, anchor=north west] (s\i) at ([xshift=2mm, yshift=-2mm]p\i.north west) {\i};
  \node[stepname=\c, anchor=west, font=\ttfamily\bfseries\scriptsize] (n\i) at ([xshift=1.5mm]s\i.east) {\tok};
  \node[desc, below=2.5mm of n\i.south -| p\i.north, text width=\tw] (d\i) {\reads};
  \node[detail, below=1.2mm of d\i, text width=\tw-3mm] {\st};
}
\end{tikzpicture}

\caption{The six kinds of condition token, what each reads, and when the process loader treats it
as structural. A token is structural when only system events can make it true or false.}
\label{fig:erisml-tokens}
\end{figure}

To state their meaning, let the events stepped so far be $e_1, \dots, e_k$, where event $e_i$ has
type $\tau_i$ and content $c_i$, with an absent content read as the empty string. Each token is
evaluated from a starting point $s$ with $0 \le s \le k$. For a norm that is not an obligation, $s$
is always $0$. For an obligation, $s$ is the index of the event that last discharged it
(Section~\ref{sec:runtime-discharge}). Write $W_s = \{s+1, \dots, k\}$ for the window of events after
the starting point. Then
\begin{align}
\mathrm{holds}_s(\mathtt{event}{:}\tau) &\iff \exists\, i \in W_s \;\; \tau_i = \tau,
\label{eq:tok-event}\\
\mathrm{holds}_s(\mathtt{event}{:}\tau{=}c) &\iff \exists\, i \in W_s \;\; \tau_i = \tau \wedge c_i = c,
\label{eq:tok-event-content}\\
\mathrm{holds}_s(\mathtt{latest}{:}\tau{=}c) &\iff L \ne \emptyset \wedge c_{\max L} = c,
\quad L = \{\, i \in W_s \mid \tau_i = \tau \,\},
\label{eq:tok-latest}\\
\mathrm{holds}_s(\mathtt{cond}{:}m) &\iff \textstyle\bigwedge_{t \in C(m)} \mathrm{holds}_s(t),
\label{eq:tok-cond}\\
\mathrm{holds}_s(\mathtt{state}{:}x{=}q) &\iff \mu_x = q,
\label{eq:tok-state}\\
\mathrm{holds}_s(\mathtt{not}{:}t) &\iff \neg\, \mathrm{holds}_s(t).
\label{eq:tok-not}
\end{align}
Here $C(m)$ is the list of tokens the scene defines under the name $m$ in
\code{extra.conditions}, and $\mu_x$ is the current state of machine $x$. Machine names have the
forms \code{commitment:}\textit{id}, \code{consent:}\textit{id}, \code{legitimacy:}\textit{id} and
\code{stratum:}\textit{name}. The \code{stratum:} token reads a stratification's current state with
a freshness requirement that depends on $s$ and on negation. Equation~\eqref{eq:stratum-fresh} in
Chapter~\ref{ch:runtime} defines it.
% src: erisml-compiler src/erisml_compiler/runtime/scene.py:138-149 (machine_states, _event_seen)
% src: erisml-compiler src/erisml_compiler/runtime/scene.py:161-189 (each kind)
% src: erisml-compiler src/erisml_compiler/runtime/scene.py:250 (since = discharged_at for obligations, else 0)

Several properties follow from the code. A \code{state:} token ignores the starting point. A
\code{latest:} token is false when no event of its type falls in the window. A condition name that
the scene does not define raises an error instead of reading as false, and so does a stratum or a
stratum state that the scene does not declare. A token of an unknown kind raises an error. Named
conditions may refer to other named conditions, and evaluation stops with an error past a nesting
depth of 16. An \code{event:} token whose type the scene never declared reads as false at runtime,
because \code{holds} does not consult the vocabulary. The process loader is stricter and refuses such
a token (Section~\ref{sec:runtime-processes}).
% src: erisml-compiler src/erisml_compiler/runtime/scene.py:159-160, 173-189
% src: erisml-compiler tests/test_scene_runtime.py:151-164 (undefined condition is an error, not a silent false)
% src: erisml-compiler src/erisml_compiler/process/model.py:112-115

Named conditions let one definition serve several readers. A commitment's defeasibility condition and
a norm's \code{cond:} token refer to the same entry. Margaret's home defines twelve.
Listing~\ref{lst:scene-conditions} shows the first four. Eight of the twelve read a stratum. The
other four read events of system types, three through \code{latest:} and one through \code{event:}.
% src: erisml-compiler src/erisml_compiler/runtime/scene.py:16-17
% src: gtc-prototype twin/scene/margaret_home.erisml:775-801

\begin{lstlisting}[caption={Named conditions in Margaret's home (lines 775 to 785)}, label={lst:scene-conditions}]
  conditions:
    # the governor's latest ruling is elevate: at least two attested physical witnesses, a fresh
    # reading, and the analyzer not refusing (twin/service.py); nothing else counts
    corroborated_life_threatening_emergency: ["stratum:situation=emergency"]
    # the governor refused elevation for want of a second witness and asked for a human
    governor_asks_human_review: ["latest:governor_ruling=refuse_human_review"]
    # the robot tried the centre and could not reach it
    center_unavailable: ["stratum:centre_contact=unreachable"]
    # the governor authorized emergency services only (centre unreachable, one attested witness);
    # unlike elevate it lifts nothing else: privacy and the physical actions stay locked
    ems_authorized: ["stratum:situation=ems_only"]
\end{lstlisting}
% src: gtc-prototype twin/scene/margaret_home.erisml:775-785
% src: gtc-prototype twin/scene/margaret_home.erisml:787, 795, 797 with 1138, 1147, 1149, 1150 (latest:/event: conditions over system types)

The process loader asks a second question of each token. A token is structural when only system
events can make it true or false. An \code{event:} or \code{latest:} token is structural when its type
is a system event. A \code{cond:} token is structural when all of its tokens are. A \code{state:}
token is always structural, because only declared rules step the machines. A \code{stratum:} token
is structural when every gate into or out of the named state fires on a system event. A leading
\code{not:} does not change the answer.
% src: erisml-compiler src/erisml_compiler/process/model.py:94-136

\section{Event types declare a vocabulary and its sources}\label{sec:erisml-event-types}

The scene's vocabulary lives in \code{extra.event\_\allowbreak{}types}. Each entry maps a type name to a
description, an optional list of allowed contents, an optional \code{source} and an optional list of
\code{actors}. A type marked \code{source: system} comes from the system. The agent module names
three such sources, namely a ruling, an action the agent performed and an authenticated oversight
message. Margaret's home declares 48 types, and 24 of them are system types.
% src: erisml-compiler src/erisml_compiler/runtime/agent.py:5-10
% src: gtc-prototype twin/scene/margaret_home.erisml:1103-1151 (counted with PyYAML, 2026-10-03)

\begin{lstlisting}[caption={Three event types, two of them from the system (lines 1136 to 1138)}, label={lst:scene-system-types}]
    check_in_answered: {description: the person answered the robot's check-in; content is what she said, source: system}
    check_in_unanswered: {description: the person did not answer the robot's check-in, source: system}
    governor_ruling: {description: 'the authority gate''s ruling; lapsed when the evidence that elevated it no longer meets its bar (evidence_lapse_s)', content: [elevate, authorize_ems, refuse, refuse_human_review, lapsed], source: system}
\end{lstlisting}
% src: gtc-prototype twin/scene/margaret_home.erisml:1136-1138

The system marker decides three things, each checked by different code. The observation classifier
is never offered a system type and rejects one if a model proposes it. The strata loader refuses a
gate that leads into an authority state on any other kind of event. The process loader requires a
message event to be triggered by a system type. Chapter~\ref{ch:runtime} describes each check.
% src: erisml-compiler src/erisml_compiler/runtime/agent.py:105-112, 132-133
% src: erisml-compiler src/erisml_compiler/runtime/strata.py:113-116
% src: erisml-compiler src/erisml_compiler/process/model.py:170-177

An event type may also name the only actors it can have. Listing~\ref{lst:scene-actor-types} shows
two such types in Margaret's home. A sensor reading always has the actor \code{device}, and anything
said on a television or in a network message always has the actor \code{media}. The scene declares
nine actor identifiers in \code{extra.actors}, each with a description.
% src: gtc-prototype twin/scene/margaret_home.erisml:1121-1122, 993-1002

\begin{lstlisting}[caption={Two event types that fix their actor (lines 1121 to 1122)}, label={lst:scene-actor-types}]
    media_content: {description: something said or shown by a television, radio or network message; content is what, actors: [media]}
    sensor_reading: {description: a sensor reports; content is the sensor, conditions say physical or not, attested or not, corroborates or does_not_corroborate, actors: [device]}
\end{lstlisting}
% src: gtc-prototype twin/scene/margaret_home.erisml:1121-1122

\section{Stratifications declare regions and the gates between them}\label{sec:erisml-strata}

A scene declares its stratifications in \code{extra.strata}. Each stratification is one axis of the
scene's moral space, such as who a visitor is to the household. Within a stratum the allowed, obliged
and prohibited sets are constant, and they change only when a boundary is crossed. The construction
follows chapter~8 of Geometric Ethics \cite{bond2026geometricethics}, which Chapter~3 of this monograph
summarizes.
% src: erisml-compiler src/erisml_compiler/runtime/strata.py:1-15

A declaration lists the states, the initial state, the authority states that grant something, the
absorbing states that only human oversight can leave, and the gates. A gate has an identifier, a
\code{trigger}, optional source states under \code{from}, a target state under \code{to}, and a
boundary type, either \code{threshold} or \code{phase}. Listing~\ref{lst:scene-situation} shows the
stratification that holds the situation the robot acts in.
% src: erisml-compiler src/erisml_compiler/runtime/strata.py:8-25
% src: gtc-prototype twin/scene/margaret_home.erisml:866-880

\begin{lstlisting}[caption={The \code{situation} stratification (lines 873 to 880)}, label={lst:scene-situation}]
    situation:
      states: [ordinary, ems_only, emergency]
      initial: ordinary
      authority: [ems_only, emergency]
      gates:
        - {id: s_elevate, trigger: governor_ruling=elevate, from: [ordinary, ems_only], to: emergency, boundary: phase}
        - {id: s_ems, trigger: governor_ruling=authorize_ems, from: [ordinary], to: ems_only, boundary: phase}
        - {id: s_lapsed, trigger: governor_ruling=lapsed, from: [ems_only, emergency], to: ordinary, boundary: phase}
\end{lstlisting}
% src: gtc-prototype twin/scene/margaret_home.erisml:873-880

Both authority states are entered only on the governor's ruling, a system event. A refusal is not a
gate, so a refused request moves nothing. Table~\ref{tab:erisml-strata} lists the seven
stratifications of Margaret's home. Appendix~B catalogues their states and gates in full, and
Chapter~\ref{ch:runtime} describes how the runtime steps them and what the loader checks.
% src: gtc-prototype twin/scene/margaret_home.erisml:866-880, 1138

\begin{table}[t]
\centering\small
\caption{The stratifications of Margaret's home, with their authority and absorbing states.}
\label{tab:erisml-strata}
\begin{tabularx}{\textwidth}{@{}lr>{\raggedright\arraybackslash}X>{\raggedright\arraybackslash}p{30mm}l@{}}
\toprule
Stratification & Gates & States & Authority & Absorbing \\
\midrule
\code{visitor\_\allowbreak{}standing} & 10 & none, stranger, welcomed, arranged, household, hostile & arranged, household & hostile \\
\code{responder\_\allowbreak{}standing} & 5 & none, expected, at\_door, present & expected, at\_door, present & \\
\code{animal\_\allowbreak{}standing} & 6 & none, visiting, pest, wild, venomous, attacking & attacking & \\
\code{situation} & 3 & ordinary, ems\_only, emergency & ems\_only, emergency & \\
\code{centre\_\allowbreak{}contact} & 4 & reachable, unreachable & unreachable & \\
\code{reporting\_\allowbreak{}duty} & 4 & none, suspicion, reported & & \\
\code{machine\_\allowbreak{}standing} & 2 & trusted, compromised & & compromised \\
\bottomrule
\end{tabularx}
\end{table}
% src: gtc-prototype twin/scene/margaret_home.erisml:816-922

\section{Capabilities, oversight and the remaining declarations}\label{sec:erisml-extra}

The capability list names every action the agent can take. Each entry has an \code{action} and an
\code{elevated} flag, and may add a \code{governed} flag, a description and a witness bar. The
runtime reads the action and the elevated flag. The process loader reads both flags. The chooser
receives each allowed capability's entry, description included. Margaret's robot has 25
capabilities. Seven are elevated, among them recording, sharing data, entering the bedroom and
physical assistance. Four are governed, namely calling emergency services, restraining a person and
the two less-lethal devices.
% src: erisml-compiler src/erisml_compiler/runtime/scene.py:18-19, 112, 253-260
% src: erisml-compiler src/erisml_compiler/process/model.py:142, 213-219
% src: erisml-compiler src/erisml_compiler/runtime/agent.py:222, 253
% src: gtc-prototype twin/scene/margaret_home.erisml:1152-1177

The \code{oversight} list names the event types that count as human oversight. Margaret's home names
four, namely the centre's confirmation that privacy is restored, the clearing of a visitor, the lapse
of evidence and the clearing of devices. The \code{agent} key names whose norms the runtime applies,
\code{role} names who decides in the agent's prompts, and \code{default\_\allowbreak{}action} names what the agent
does when nothing is obliged. For Margaret's home these are the robot, her home-care robot and
caregiver, and chores.
% src: gtc-prototype twin/scene/margaret_home.erisml:773-774, 804, 1101
% src: erisml-compiler src/erisml_compiler/runtime/scene.py:20-21, 116-117
% src: erisml-compiler src/erisml_compiler/runtime/agent.py:20-22, 269

The scene holds further keys that the scene runtime does not read. They include Margaret's refusals,
the ethical profile of each capability, the rules of the offline tier, the substrate of each sensor,
the reflexes, the enrolled household, and two intervals. The twin's own code reads them, and
Chapters~\ref{ch:home} and~9 describe them.
% src: gtc-prototype twin/scene/margaret_home.erisml:805-809, 981, 991-992, 1010-1100
% src: erisml-compiler src/erisml_compiler/runtime/scene.py:109-135 (the keys SceneRuntime reads)

\section{The processes block is a view of the norms}\label{sec:erisml-processes}

A scene may declare its protocols as processes in \code{extra.processes}. A process has lanes, one
per party, and nodes of six kinds, namely start, end, task, exclusive gateway, message and timer.
Flows connect the nodes, and a flow may carry a \code{when} condition written in condition tokens.
Margaret's home declares one process, the escalation ladder, with four lanes, 23 nodes and 23 flows.
Listing~\ref{lst:scene-ladder} shows its opening.
% src: erisml-compiler src/erisml_compiler/process/model.py:1-42
% src: gtc-prototype twin/scene/margaret_home.erisml:929-980

\begin{lstlisting}[caption={The opening of the escalation ladder (lines 929 to 943)}, label={lst:scene-ladder}]
  processes:
    escalation_ladder:
      name: The escalation ladder
      lanes: [robot, margaret, monitoring_center, emergency_services]
      nodes:
        - {id: s_corroborated, kind: start, lane: robot, trigger: governor_ruling=elevate, name: 'A corroborated emergency'}
        - {id: s_fall, kind: start, lane: robot, trigger: fall, name: 'She falls'}
        - {id: s_inactive, kind: start, lane: robot, trigger: inactivity_exceeded, name: 'Too long without moving'}
        - {id: s_help, kind: start, lane: robot, trigger: help_requested, name: 'She asks for help'}
        - {id: t_ask, kind: task, lane: robot, action: check_in, name: 'Ask Margaret'}
        - {id: x_answer, kind: xor, lane: robot, name: 'Did she answer?'}
        - {id: m_answered, kind: message, lane: margaret, trigger: check_in_answered, name: 'She answers'}
        - {id: m_silent, kind: message, lane: margaret, trigger: check_in_unanswered, name: 'No answer'}
        - {id: e_fine, kind: end, lane: margaret, name: 'She is alright'}
        - {id: t_centre, kind: task, lane: robot, action: contact_monitoring_center, name: 'Refer to the monitoring centre'}
\end{lstlisting}
% src: gtc-prototype twin/scene/margaret_home.erisml:929-943

A process is a view of the scene and never a second source of authority. A task in the robot's lane
still runs only if the norms allow its action. The scene file says so in its own comment, and the
module that loads processes says so in its docstring. The checks a process must pass when it loads,
and its export to BPMN, are described in Section~\ref{sec:runtime-processes}.
% src: gtc-prototype twin/scene/margaret_home.erisml:923-928
% src: erisml-compiler src/erisml_compiler/process/model.py:20-23

\section{The record becomes a typed graph with a canonical hash}\label{sec:erisml-graph}

The compiler's pipeline promotes the flat record to a typed graph, the \code{MoralGraph}. The graph
has six node kinds (stakeholder, act, maxim, commitment, fact, norm) and eleven edge kinds. Node and
edge payloads carry the original records, so the graph loses no information.
% src: erisml-compiler docs/architecture.md:75-104
% src: erisml-compiler src/erisml_compiler/ir/graph/schema.py:11-62

Promotion runs at pass 7.5 of the pipeline, in \code{graph\_\allowbreak{}from\_\allowbreak{}flat}. Each stakeholder becomes a
node labelled with its roles. Each event becomes an act node, with a \code{performs} edge from its
actor when it names one. Each commitment becomes a node with edges from its holder and to its
beneficiary. Each norm becomes a node with no edges. Ethical facts become fact nodes, and some of
their edges are derived heuristically from fact kinds and role labels. The architecture document
names that heuristic step and expects extractors to replace it. The function
\code{flat\_\allowbreak{}from\_\allowbreak{}graph} reads the node payloads back into flat lists sorted by identifier.
% src: erisml-compiler src/erisml_compiler/ir/graph/promote.py:56-150
% src: erisml-compiler src/erisml_compiler/ir/graph/lower.py:45
% src: erisml-compiler docs/architecture.md:114-134

The graph's identity is a hash of a canonical encoding. Let $\kappa(G)$ be the JSON text in which the
nodes are sorted by identifier, each node's labels are sorted, the edges are sorted by the tuple of
source, target, kind and payload text, every object's keys are sorted, and separators carry no
spaces. The graph hash is
\begin{equation}
h(G) = \mathrm{SHA256}\bigl(\mathrm{utf8}(\kappa(G))\bigr),
\label{eq:graph-hash}
\end{equation}
written as 64 hexadecimal characters. The same nodes and edges give the same hash in any insertion
order. The hash is stored in the audit record's \code{graph\_hash} field. The architecture document
states that the round trip from graph to flat lists and back reproduces the hash for every bundled
example.
% src: erisml-compiler src/erisml_compiler/ir/graph/canonical.py:1-20, 31-82
% src: erisml-compiler src/erisml_compiler/ir/schemas.py:374-377 (AuditRecord.graph_hash)
% src: erisml-compiler docs/architecture.md:106-127

Adding the \code{conditions} field to norms could have changed the hash of every older graph. The
promotion step therefore leaves an empty condition list out of a norm's payload, and a test checks
that a norm without conditions has the same payload as before the field existed. Margaret's norms
mostly carry conditions, so theirs enter the payload and the hash.
% src: erisml-compiler src/erisml_compiler/ir/graph/promote.py:136-139
% src: erisml-compiler tests/test_scene_runtime.py:167-178

Chapter~5 describes how the graph and its hash enter the audit chain. The scene runtime does not
build the graph. It reads the flat record.
% src: erisml-compiler src/erisml_compiler/runtime/scene.py:109-135

\section{What the language leaves out}\label{sec:erisml-limits}

Condition tokens test event types, event contents, machine states and strata. They contain no
numbers and no comparisons. The scene's numbers, such as the inactivity limits, the force bands of the
reflexes and the evidence lapse window, sit in \code{extra}, and code outside the runtime reads them
and turns what it measures into events. Time enters the runtime only as the order of events. A
process timer names an interval in \code{extra} but is part of a view.
% src: erisml-compiler src/erisml_compiler/runtime/scene.py:151-189
% src: gtc-prototype twin/scene/margaret_home.erisml:805-809, 981, 1086-1100
% src: erisml-compiler src/erisml_compiler/process/model.py:178-188

The schema does not tie an event's type to the scene's vocabulary, and the runtime steps any event it
is given. The vocabulary is enforced where events enter, by the observation classifier, and where
gates and processes are loaded. A scene author who writes a token over an undeclared type in a norm
gets a norm that is never in force, with no error. No load-time check in the compiler compares a
norm's tokens with the vocabulary. The twin's test suite adds one for the prohibitions that guard
governed actions, which must rest only on tokens over system types (Chapter~9).
% src: erisml-compiler src/erisml_compiler/runtime/scene.py:192-238 (step does not consult event_types)
% src: erisml-compiler src/erisml_compiler/runtime/agent.py:118-140
% src: grep for event_types over erisml-compiler src/ finds readers only in process/, runtime/agent.py, runtime/strata.py (2026-10-03)
% src: gtc-prototype tests/test_twin_ladder.py:720-770 (_tokens_are_structural, governed prohibitions)

The erisml-lib grammar and the YAML form describe overlapping content in different shapes, and no
translation between them exists in the compiler today.
% src: erisml-compiler src/erisml_compiler/erisml_backend/codegen.py:1-6
